% =====================================================================
%  MIC Summer-of-Code  --  Week 4 (supplementary)
%  Kernel Methods: the kernel trick, RKHS, SVMs, and Kernel PCA
%
%  Self-contained handbook chapter.  Compile with pdflatex (twice) or Overleaf.
%      pdflatex MIC_Kernel_Methods.tex
%      pdflatex MIC_Kernel_Methods.tex
%
%  Built from Prof. Suyash P. Awate's CS-736 KernelMethods deck, with
%  attribution, for the use of mentees on this project.
% =====================================================================
\documentclass[11pt]{article}

\usepackage[a4paper,margin=1in]{geometry}
\usepackage{amsmath,amssymb,amsthm,mathtools,bm}
\usepackage{enumitem}
\usepackage{booktabs}
\usepackage{xcolor}
\usepackage{microtype}
\usepackage[most]{tcolorbox}
\usepackage{titlesec}
\usepackage[colorlinks=true,linkcolor=blue!45!black,urlcolor=blue!55!black,
            citecolor=blue!45!black]{hyperref}

\titleformat{\section}{\Large\bfseries\sffamily\color{blue!40!black}}{\thesection}{0.6em}{}
\titleformat{\subsection}{\large\bfseries\sffamily\color{blue!32!black}}{\thesubsection}{0.5em}{}

\newtcolorbox{keybox}[1][]{colback=blue!4!white,colframe=blue!40!black,
  fonttitle=\bfseries\sffamily,title=#1,boxrule=0.6pt,arc=2pt,left=4pt,right=4pt,top=3pt,bottom=3pt}
\newtcolorbox{notebox}[1][Aside]{colback=black!3!white,colframe=black!45!white,
  fonttitle=\bfseries\sffamily,title=#1,boxrule=0.5pt,arc=2pt,left=4pt,right=4pt,top=3pt,bottom=3pt}

\newcommand{\R}{\mathbb{R}}
\newcommand{\norm}[1]{\left\lVert #1 \right\rVert}
\newcommand{\inner}[2]{\langle #1,\,#2\rangle}
\newcommand{\deck}[1]{\textsf{\small[#1]}}
\DeclareMathOperator*{\argmin}{arg\,min}

\title{\textbf{Kernel Methods}\\[2pt]
\large The kernel trick, RKHS, SVMs, and Kernel PCA\\[6pt]
\normalsize Medical Image Computing -- Summer of Code -- Week 4 (supplementary)}
\author{}
\date{}

\begin{document}
\maketitle
\vspace{-2.2em}

\begin{keybox}[How to read this chapter]
This is the supplementary mathematical capstone of Week 4. Its one big idea is
beautiful and reusable: when your data is too tangled for a \emph{linear} method,
secretly map it into a much higher-dimensional space where it \emph{becomes}
linear --- and do it without ever paying for that space, using a \textbf{kernel}.
The same trick upgrades clustering (kernel k-means), dimensionality reduction
(\textbf{kernel PCA}), classification (\textbf{SVMs}), regression (kernel ridge),
and the shape models of this week's main chapter. Keep the \deck{KernelMethods}
deck open; do the kernel-PCA notebook, where linear PCA fails on two moons and a
kernel pulls them apart. You will also recognise old friends: RKHS, positive
definiteness, and Gram matrices all came up in Weeks 1--4.
\end{keybox}

\tableofcontents
\bigskip

% =====================================================================
\section{When linear methods fail}
% =====================================================================

Two complementary modelling styles run through machine learning
\deck{KernelMethods, slides 3--4}: \textbf{generative} (model each class's
distribution, e.g.\ a Gaussian per cluster, and assign a test point to the
highest-density class) and \textbf{discriminative} (model the boundary directly,
e.g.\ a separating hyperplane as in an SVM). Both have a linear comfort zone, and
both break outside it \deck{KernelMethods, slide 5}: a discriminative hyperplane
cannot separate classes whose true boundary is a curve, and a single Gaussian
cannot model a cluster that is a ring or a crescent rather than an ellipsoid.
Kernel methods are how we keep the simple linear machinery while handling these
nonlinear cases.

% =====================================================================
\section{The feature-map idea}
% =====================================================================

The move is to \textbf{map the data into a higher-dimensional space} where the
nonlinear problem becomes linear \deck{KernelMethods, slides 6--7}. The classic
picture: two classes arranged as concentric rings in 2D cannot be split by any
line, but lift each point $[x,y]$ to $[x,\,y,\,x^2+y^2]$ and the inner ring rises
to a different height than the outer one --- now a flat plane separates them
cleanly. More generally, choose a feature map $\phi(\cdot)$ so that, in the mapped
space, the clusters are ellipsoidal (good for Gaussians) or linearly separable
(good for hyperplanes). The catch: a map powerful enough to do this is usually
\emph{very} high- (even infinite-) dimensional, so computing and storing
$\phi(x)$ explicitly is hopeless.

% =====================================================================
\section{The kernel trick}
% =====================================================================

Here is the escape \deck{KernelMethods, slide 8}. Almost every linear algorithm
(PCA, SVM, ridge regression, k-means) can be rewritten so that data appears
\emph{only} through \textbf{inner products} $\inner{x}{y}$. So if we have a
function that returns the inner product \emph{in the mapped space} without
building $\phi$,
\[
k(x,y) = \inner{\phi(x)}{\phi(y)},
\]
we can run the whole algorithm on $k$ and never touch $\phi(x)$. That function is
the \textbf{kernel}, and using it in place of explicit features is the
\textbf{kernel trick}. The identity map gives the ordinary dot product
$k(x,y)=\inner{x}{y}$; the workhorse nonlinear choice is the \textbf{Gaussian} /
\textbf{radial basis function (RBF)} kernel
\[
k(x,y) = \exp\!\big(-a\,\norm{x-y}^2\big),
\]
whose implicit feature space is \emph{infinite}-dimensional \deck{KernelMethods,
slide 9}. In practice you \emph{choose the kernel first} and never write $\phi$
down at all.

\begin{notebox}[Worked example: the RBF kernel hides an infinite feature map]
Expand a polynomial kernel $k(x,y)=(\inner{x}{y})^2$ in 2D: it equals
$\inner{\phi(x)}{\phi(y)}$ for $\phi(x)=(x_1^2,\sqrt2\,x_1x_2,x_2^2)$ --- a
\emph{finite} lift you could write out. The RBF kernel is the limit of this idea
to all polynomial degrees at once: its $\phi$ has infinitely many coordinates, so
you truly \emph{cannot} write it down --- yet $k(x,y)=e^{-a\norm{x-y}^2}$ is a
one-line computation. That gap (a one-line kernel standing in for an infinite
feature vector) is the whole payoff.
\end{notebox}

% =====================================================================
\section{What makes a kernel valid: inner products, Hilbert space, RKHS}
% =====================================================================

Not every two-argument function is a legal kernel; it must be the inner product
of \emph{some} feature map. The structure that guarantees this is a
\textbf{Hilbert space} --- a vector space equipped with an inner product
\deck{KernelMethods, slides 10--12}. An inner product must satisfy three axioms:
\begin{enumerate}[nosep,leftmargin=*]
  \item \textbf{conjugate symmetry} $\inner{f}{g}=\overline{\inner{g}{f}}$;
  \item \textbf{linearity} in the first argument;
  \item \textbf{positive definiteness} $\inner{f}{f}\ge0$, with equality only for
  the zero element.
\end{enumerate}
From an inner product you get lengths and angles (and hence distances), which is
all any of our geometric algorithms need.

The precise object is a \textbf{Reproducing Kernel Hilbert Space (RKHS)}
\deck{KernelMethods, slides 14--16}: a Hilbert space of functions in which
evaluation is continuous, represented by a \textbf{reproducing kernel}
$k(x,\cdot)$ with the magic \emph{reproducing property} $\inner{k(x,\cdot)}{k(y,\cdot)}=k(x,y)$.
The chain of facts to remember:
\begin{itemize}[nosep,leftmargin=*]
  \item every reproducing kernel \emph{is} a kernel (take $\phi(x)=k(x,\cdot)$);
  \item every kernel is a \textbf{positive-definite (PD)} function, because
  $\inner{\sum_i a_i\phi(x_i)}{\sum_i a_i\phi(x_i)}=\norm{\sum_i a_i\phi(x_i)}^2\ge0$;
  \item conversely (the \textbf{Moore--Aronszajn} theorem, the kernel analogue of
  Mercer's theorem), \emph{every} symmetric PD function is the reproducing kernel
  of a unique RKHS.
\end{itemize}
So the practical test ``is my $k$ a valid kernel?'' becomes ``is $k$ symmetric and
positive-definite?'' --- and PD-ness is the same property you met for covariance
and Gram matrices in earlier weeks.

% =====================================================================
\section{Support Vector Machines (briefly)}
% =====================================================================

The most famous kernel method is the \textbf{Support Vector Machine}
\deck{KernelMethods, slide 4}. For linearly separable classes the SVM finds the
\textbf{maximum-margin} hyperplane --- the separator with the widest empty gap on
both sides --- which generalises better than an arbitrary separator. Solving it is
a constrained convex optimisation whose \textbf{KKT conditions} (the
inequality-constrained optimality conditions from the segmentation chapter)
reveal that the solution depends only on a few boundary points, the
\textbf{support vectors}. Crucially, the SVM's data appears only through inner
products, so \emph{kernelising} it --- replace $\inner{x}{y}$ by $k(x,y)$ --- yields
a nonlinear classifier (e.g.\ an RBF-SVM with curved boundaries) for free. SVMs
were the dominant classifier before deep learning and remain excellent on small,
tabular, or feature-engineered medical data.

% =====================================================================
\section{Kernel PCA}
% =====================================================================

Recall \textbf{PCA}: centre the data, form the sample covariance, and take its
top eigenvectors as the directions of greatest variation \deck{KernelMethods,
slide 21}. It is perfect for ellipsoidal (Gaussian) data and poor for anything
curved --- a ring's principal axes tell you nothing useful \deck{KernelMethods,
slide 22}. \textbf{Kernel PCA} does PCA in the mapped space, capturing
\emph{nonlinear} structure \deck{KernelMethods, slides 23--27}.

The trick is that each eigenvector $V$ of the covariance in feature space must lie
in the \emph{span of the mapped data}, $V=\sum_i \alpha_i\phi(x_i)$, so we only
need the coefficients $\alpha$. Substituting reduces the whole eigenproblem to one
on the $N\times N$ \textbf{Gram matrix} $G_{ij}=\inner{\phi(x_i)}{\phi(x_j)}=k(x_i,x_j)$:
\[
M\lambda\,\alpha = G\,\alpha,
\]
i.e.\ the kernel-PCA components are the eigenvectors of the (centred) Gram matrix
--- computed entirely from kernel evaluations, never from $\phi$. A new point is
projected by the ``feature'' $\inner{\phi(x)}{V}=\sum_i\alpha_i\,k(x_i,x)$. Plotting
these features over the input space shows \emph{nonlinear} modes of variation that
adapt to the data's shape, and projecting onto the top kernel-PCA components
separates non-ellipsoidal clusters that linear PCA cannot \deck{KernelMethods,
slides 25--27}.

\begin{notebox}[Why ``eigenvector lies in the span of the data'' works]
The covariance $C=\frac1M\sum_i\phi(x_i)\phi(x_i)^{\!\top}$ acts on any vector by
taking inner products with the $\phi(x_i)$, so $CV$ is automatically a combination
of the $\phi(x_i)$. Hence any eigenvector with non-zero eigenvalue is too --- it
lives in the (at most $M$-dimensional) span of the data, even if the feature space
is infinite-dimensional. That single observation is what collapses an
infinite-dimensional eigenproblem to an $M\times M$ one, the \textbf{representer
theorem} in miniature.
\end{notebox}

\begin{notebox}[The pre-image problem]
Kernel PCA gives you a point \emph{in feature space} (a combination of
$\phi(x_i)$'s). To \emph{see} it as an image you need its \textbf{pre-image} ---
an input $x$ whose $\phi(x)$ is closest to that feature-space point. There is
usually no exact pre-image (the feature point may not be $\phi$ of anything), so
it is found by optimisation. This is exactly the step assignment \texttt{a4} asks
for when denoising shape segmentations with kernel PCA.
\end{notebox}

% =====================================================================
\section{A unifying view, and kernel ridge regression}
% =====================================================================

Kernels also connect to density estimation: the squared distance from a mapped
point $\phi(z)$ to the mean of mapped data, for an RBF kernel, is (up to
constants) a \textbf{Gaussian mixture / kernel density estimate} in the input
space \deck{KernelMethods, slides 19--20}. So ``close to the mean in RKHS'' means
``high data density in input space'' --- kernel methods and the GMMs of the
segmentation chapter are two views of the same object.

Finally, \textbf{kernel ridge regression} kernelises the Tikhonov-regularised
least squares of Week 3 \deck{KernelMethods, slides 30--32}: a matrix identity
turns the high-dimensional ridge solution into one expressed through the Gram
matrix, and the prediction at a new point is $\sum_i\alpha_i\,k(x_i,x)$ --- again
all kernel, no $\phi$. This is the \textbf{representer theorem} at work: the
optimal function is a weighted sum of kernels centred on the data.

\begin{keybox}[The kernel recipe, in one line]
Take any linear algorithm written with inner products $\inner{x}{y}$, replace
each by a kernel $k(x,y)$, and you get its nonlinear version --- PCA $\to$ kernel
PCA, ridge $\to$ kernel ridge, k-means $\to$ kernel k-means, the SVM's nonlinear
form --- all without ever building the feature space.
\end{keybox}

% =====================================================================
\section{Kernel methods in medical image computing}
% =====================================================================

Kernel methods appear across the course's tasks \deck{KernelMethods, slide 29}:
kernel \textbf{k-means / FCM / GMM} for segmentation of non-ellipsoidal intensity
clusters; kernel \textbf{PCA} for nonlinear \emph{statistical shape models}
(assignment \texttt{a4}) and dimensionality reduction; kernel \textbf{ridge
regression} for image restoration (denoising, motion correction); and kernel
classifiers for fMRI and brain-connectivity analysis.

\begin{notebox}[Common pitfalls]
\begin{itemize}[nosep,leftmargin=*]
  \item \textbf{Kernel \& bandwidth choice.} The RBF width $a$ controls
  everything; too wide is near-linear, too narrow overfits. Tune it (honestly
  --- ``grad-student descent'' is the deck's joke for fiddling until it works).
  \item \textbf{Feature scaling.} RBF uses $\norm{x-y}$, so standardise features
  first or one dimension dominates.
  \item \textbf{It scales as $N^2$.} The Gram matrix is $N\times N$; kernel methods
  get expensive for large datasets (a reason deep nets took over at scale).
  \item \textbf{Pre-image.} You cannot always map a feature-space result back to an
  image exactly; it needs an optimisation.
\end{itemize}
\end{notebox}

\begin{keybox}[Do the notebook]
\texttt{notebooks/03\_kernel\_pca\_toy.ipynb} (NumPy + Matplotlib) builds two
interlocking ``moons'', shows that \textbf{linear PCA} cannot separate them, then
computes \textbf{kernel PCA} with an RBF kernel (centred Gram matrix
eigen-decomposition) and shows the classes falling cleanly apart in the top two
kernel components --- the kernel trick, made visible.
\end{keybox}

% =====================================================================
\section*{Resources (watch one, read one)}
\addcontentsline{toc}{section}{Resources}
% =====================================================================

\textbf{Kernels \& RKHS}
\begin{itemize}[nosep,leftmargin=*]
  \item C.\ Shalizi, \href{https://www.stat.cmu.edu/~cshalizi/350/lectures/16/lecture-16.pdf}{\emph{Mercer's theorem \& RKHS}} (concise notes).
  \item Sch\"olkopf \& Smola, \emph{Learning with Kernels}, ch.~1--2 (the standard reference).
\end{itemize}

\textbf{SVMs}
\begin{itemize}[nosep,leftmargin=*]
  \item MIT 6.034, \href{https://www.youtube.com/watch?v=_PwhiWxHK8o}{\emph{Support Vector Machines}} (Patrick Winston lecture).
  \item Boyd \& Vandenberghe, \href{https://web.stanford.edu/~boyd/cvxbook/}{\emph{Convex Optimization}}, ch.~5 (KKT conditions).
\end{itemize}

\textbf{Kernel PCA}
\begin{itemize}[nosep,leftmargin=*]
  \item scikit-learn, \href{https://scikit-learn.org/stable/auto_examples/decomposition/plot_kernel_pca.html}{kernel PCA example}.
  \item Sch\"olkopf, Smola \& M\"uller (1998), \emph{Nonlinear Component Analysis as a Kernel Eigenvalue Problem} (the kernel-PCA paper).
\end{itemize}

\vfill
\begin{center}\small\itshape
Built for the MIC Summer-of-Code from Prof.\ Suyash P.\ Awate's CS-736
KernelMethods deck, with attribution. For mentee use on this project; please do
not redistribute the bundled decks outside it.
\end{center}

\end{document}
